\documentclass[a4paper]{article}

% if you need to pass options to natbib, use, e.g.:
%     \PassOptionsToPackage{numbers, compress}{natbib}
% before loading neurips_2022


% ready for submission
%\usepackage{ofu_xai_2022}

\input{math_commands.tex}

% to compile a preprint version, e.g., for submission to arXiv, add add the
% [preprint] option:
%     \usepackage[preprint]{ofu_xai_2022}


% to compile a camera-ready version, add the [final] option, e.g.:
\usepackage[final, nonatbib]{ofu_xai_2022}

% to avoid loading the natbib package, add option nonatbib:
%    \usepackage[nonatbib]{ofu_xai_2022}


\usepackage[utf8]{inputenc} % allow utf-8 input
\usepackage[T1]{fontenc}    % use 8-bit T1 fonts
\usepackage{hyperref}       % hyperlinks
\usepackage{url}            % simple URL typesetting
\usepackage{booktabs}       % professional-quality tables
\usepackage{amsfonts}       % blackboard math symbols
\usepackage{nicefrac}       % compact symbols for 1/2, etc.
\usepackage{microtype}      % microtypography
\usepackage{xcolor}         % colors
\usepackage{longtable}
\usepackage{dsfont}

\usepackage[numbers]{natbib}
%%%%

\title{Penultimate-Layer Feature-Based OOD Detection: Feature Shaping and Its Limits}




\author{%
  Martin Simons\\
  Otto-Friedrich University of Bamberg\\
  Aufbaustr. 11, 96049 Bamberg, Germany\\
  \texttt{martin.simons@stud.uni-bamberg.de} \\[0.5cm]
  Seminar: xAI-Sem-M \\
  Degree: M.Sc. (CitH) \\
  Matriculation \#: 2188932
}


\begin{document}


\maketitle
\def\va{{\bm{a}}}

\begin{abstract}
  Post-hoc out-of-distribution (OOD) detection equips a frozen, pre-trained
  classifier to flag inputs unlike its training data, without retraining. We examine
  five feature-based methods that read the penultimate-layer representation: four
  feature-shaping rules---ReAct, ASH, VRA, and Optimal-FS---that transform the feature
  vector before rescoring, and BLOOD, which instead measures the smoothness of the
  network's between-layer transformation. Although the methods share a layer and, in
  their mathematical form, assume no particular architecture, they were developed on
  different architecture families and never co-evaluated under one protocol. Reading
  them together, we find three things: the shaping rules succeed only on the
  architecture they were designed on, with Optimal-FS and ReAct transferring while pruning rules collapse; the
  cluster unifies only partially, as whole-vector ASH and the Jacobian-based BLOOD lie
  outside the optimization that explains the element-wise rules. BLOOD, used as a
  control, shows the OOD signal is not carried by the feature values alone.
\end{abstract}


\section{Introduction}
\label{intro}

% Introduction --- motivation, the post-hoc feature-based setting, the shaping
% cluster and BLOOD, and what this report does with them.

Out-of-distribution (OOD) detection---recognizing when an input falls outside the
distribution a model was trained on---has become a basic requirement for deploying
machine-learning systems in the open world, where the inputs met at test time cannot
all be anticipated during training and a confidently wrong prediction on an unfamiliar
input can be costly. The need is sharpened by a shift in how models are built:
real world applications increasingly rely on large pre-trained models that were not fine-tuned and
cannot easily be retrained, used as frozen checkpoints. This makes \emph{post-hoc}
OOD detection, methods that attach to an already-trained network at inference time,
add no training cost, and leave its weights untouched \cite{survey}, the dominant practical setting,
and the one this report studies.

Within this setting, \emph{feature-based} detectors read intermediate
representations rather than the output logits alone --- generally features of the penultimate layer.
 Most form a tight family we call the \emph{ feature shaping cluster}: ReAct~\citep{react} (rectified activations),
ASH~\citep{ash} (activation shaping), VRA~\citep{vra} (variational rectified
activation), and Optimal-FS~\citep{optimalfs}\footnote{Since \citet{optimalfs} do not assign a shorthand to their presented method, we have assigned their method the name \emph{Optimal-FS} after the paper's title \emph{Towards Optimal Feature-Shaping Methods} as a reference for this report.} each insert a fixed transformation that rewrites the feature
vector before the logits are recomputed, sharpening the separation between
in-distribution (ID) and OOD inputs while leaving ID classification essentially
unchanged. BLOOD~\citep{blood} (between-layer out-of-distribution detection) offers an
alternative view of the same representation: rather than reshaping the feature vector,
it measures how smoothly the network transforms it, on the premise that a trained model
processes ID inputs more smoothly than OOD ones. Nothing in the mathematical setup of
these methods refers to a particular architecture: each is defined for any frozen
classifier with a penultimate layer. On paper, none of the methods are restricted to any model family.

This report reads the five methods together. The shaping cluster rests on a shared
assumption, that the discriminative OOD signal resides in the penultimate feature
vector and can be recovered by reshaping it. BLOOD, by reading a property of the
transformation rather than editing the vector, serves as a control on that assumption.
Because the methods were developed on different architecture families and never
co-evaluated under a single protocol, we account for this by asking only three questions: why the same method succeeds
on some architectures and fails on others, how far the shaping cluster can be unified,
and whether the OOD signal lies in the feature values or in the transformation applied
to them.

\section{Related Work}
\label{related}


\subsection{Positioning within the Survey Taxonomy}
\label{subsec:taxonomy}

We situate the methods in the mechanism-oriented taxonomy of \citet{survey}. Its
first cut separates \emph{training-driven} methods, which alter the objective or the
training data, from \emph{training-agnostic} ones that operate on a frozen model; a
second cut separates \emph{test-time adaptive} methods, which update running
statistics from test batches, from \emph{post-hoc} methods, which score each input
from a single forward pass. All five methods studied here are training-agnostic and
post-hoc.

Post-hoc detectors differ in which model-internal quantity they read: the logits
(e.g.\ MSP\cite{MSP}, MaxLogit, Energy\cite{energy}), a feature-space distance (e.g.\ Mahalanobis \cite{mahalanobis}), a
gradient, the penultimate representation, or an estimated density. This report
concentrates on the \emph{feature-based} family, which reads the penultimate
activation vector $\mathbf{z} \in \mathbb{R}^{M}$ produced before the linear head
$\mathbf{W} \in \mathbb{R}^{M \times C}$. Within it, \citet{survey} separate methods
that reshape $\mathbf{z}$ directly from those that read a statistical property of it.
The first branch \emph{feature shaping}, replaces $\mathbf{z}$ with $ \bar{\mathbf{z}}=f(\mathbf{z})$
before the logits are recomputed, with $f$ chosen to sharpen the ID/OOD gap while
leaving ID accuracy intact; ReAct~\cite{react}, VRA~\cite{vra}, and
Optimal-FS~\cite{optimalfs} form a progression from hand-designed rules toward
Optimal-FS's optimization view, which recovers the earlier rules as approximations to
a single ID-only optimum as an element-wise operation, complemented by ASH~\cite{ash}, which performs whole vector modifications. BLOOD~\cite{blood} sits in the second branch: it leaves
$\mathbf{z}$ untouched and measures the smoothness of the network's intermediate
transformation through the feature Jacobian, and here serves as the control that
isolates the shaping cluster's shared assumption: that the OOD signal resides in
$\mathbf{z}$ itself.

% Follows \ref{subsec:taxonomy} within Related Work.
\subsection{A Second Axis: Access Requirements}
\label{subsec:access}

Beyond \emph{which} quantity a method reads, \citet{blood} note a second,
deployment-relevant axis, the \emph{access} a method requires. \emph{Black-box}
methods need only inputs and outputs; \emph{white-box} methods also need the weights
but no data; \emph{open-box} methods additionally need training data, an OOD
validation set, or training intervention. This taxonomy further differentiates the analyzed methods:
ASH~\cite{ash} and BLOOD~\cite{blood} use no training data and are white-box, whereas
ReAct~\cite{react}, VRA~\cite{vra}, and Optimal-FS~\cite{optimalfs} fit thresholds or
shaping functions to ID data and are open-box. The strongest open-box reference point, the Mahalanobis
distance detector~\cite{mahalanobis}---which fits a class-conditional Gaussian to ID
features and scores by distance to the nearest class mean---is used as a benchmark against
all methods; it appears in the Optimal-FS and BLOOD experiments, to which we defer it.

\section{Methodology}
\label{sec:methodology}

We consider a classifier pre-trained on in-distribution data and held frozen.
For an input $\mathbf{x}$, the backbone produces a penultimate feature vector
$\mathbf{z} \in \mathbb{R}^{M}$, which a linear head $\mathbf{W}$ maps to $C$ class
logits. A post-hoc detector reduces the network's response to a scalar score and
flags the sample once that score crosses a
threshold. Detection is thus a thresholding rule
\begin{equation}
G_\lambda : \mathbb{R} \to \{\text{ID}, \text{OOD}\},
\end{equation}
parameterized by a threshold $\lambda$, that labels $\mathbf{x}$ by comparing its score
to $\lambda$. The rule $G_\lambda$ is common to every method; the methods differ in how
that score is \emph{produced} from the frozen network and in its orientation (which
side of $\lambda$ denotes ID). Most produce it by applying a logit score
$s : \mathbb{R}^{C} \to \mathbb{R}$ (MSP, MaxLogit, or energy) to the network's logits;
BLOOD is the exception, skipping $s$ and computing its scalar directly. Detection quality is reported
with the two standard metrics, AUROC (area under the receiver operating characteristic
curve, a threshold-free measure of separability) and FPR95 (false positive rate at
$95\%$ true positive rate). Fixing $\lambda$ is orthogonal to the score and identical
across methods---AUROC is threshold-free, and FPR95 fixes $\lambda$ at $95\%$
true-positive rate for every method alike---so the white-box/open-box distinction
(Section~\ref{subsec:access}) concerns only what computing the score requires, not
threshold calibration.

\subsection{Feature Shaping}
\label{subsec:shaping}

Feature shaping inserts a transformation $f : \mathbb{R}^{M} \to \mathbb{R}^{M}$ at
the penultimate layer, reshaping $\mathbf{z}$ before the logits are recomputed and
scored. The detector then factors into a shaping stage and a scoring stage,
\begin{equation}
f : \mathbb{R}^{M} \to \mathbb{R}^{M}, \qquad
s : \mathbb{R}^{C} \to \mathbb{R}, \qquad
\mathbf{x} \mapsto s\big(\mathbf{W} f(\mathbf{z})\big),
\end{equation}
and the sample is flagged by $G_\lambda\big(s(\mathbf{W} f(\mathbf{z}))\big)$.
The score $s$ is a standard logit-based score (MSP, MaxLogit, or energy), for which
a lower value indicates OOD; crucially, the choice of $s$ is orthogonal to the
shaping function $f$, a separation Optimal-FS later exploits. Setting
$f = \mathrm{id}$ recovers the unshaped baseline. Shaping functions are
\emph{element-wise} when $f$ acts coordinatewise, $\bar{z}_i = f(z_i)$, and
\emph{whole-vector} when the transformation of one coordinate depends on the others.

\subsubsection{ReAct and VRA}
\label{subsec:react-vra}

ReAct~\citep{react} and VRA~\citep{vra} are both element-wise. ReAct clips each
activation at an upper threshold $c$,
\begin{equation}
f_{\mathrm{ReAct}}(z_i) = \min(z_i, c),
\end{equation}
with $c$ set to a high percentile (typically the $90$th) of the activations
estimated on ID data. The motivation is that OOD inputs tend to produce a few
abnormally large activations; capping them suppresses the OOD score more than the ID
score while leaving most ID activations, which fall below $c$, untouched.

VRA starts from a variational analysis of which operation maximally separates ID and
OOD activations, concluding that the ideal rule suppresses both abnormally low and
abnormally high activations. Its canonical piecewise realization, VRA-P, uses two
ID-estimated quantile thresholds $\alpha < \beta$,
\begin{equation}
f_{\mathrm{VRA\text{-}P}}(z_i) =
\begin{cases}
0 & z_i < \alpha, \\
z_i & \alpha \leq z_i \leq \beta, \\
\beta & z_i > \beta .
\end{cases}
\end{equation}
This generalizes ReAct, recovered as the special case $\alpha = 0$ \cite{vra}: VRA-P
adds a lower cut that zeroes the abnormally small activations ReAct leaves in place.
Both methods fit their thresholds to ID data and are therefore data-dependent,
placing them in the open-box category (cf.\ Section~\ref{subsec:access}).

\subsubsection{ASH}
\label{subsec:ash}

ASH~\citep{ash} operates on the whole vector. For each sample it computes the $p$th
percentile $t$ of that sample's own activations, prunes everything below $t$ to
zero, and then treats the surviving activations in one of three ways: ASH-P leaves
them unchanged, ASH-B sets them all to a common constant, and ASH-S---the version we
take as canonical---scales the survivors by a factor that restores the pre-pruning
activation mass,
\begin{equation}
\label{eq:ash}
[f_{\mathrm{ASH\text{-}S}}(z)]_{i} =
\begin{cases}
z_i \cdot \exp\!\left( s_1 / s_2 \right) & z_i \geq t,\\[2pt]
0 & z_i < t,
\end{cases}
\qquad
s_1 = \sum_j z_j, \quad s_2 = \sum_{j : z_j \geq t} z_j ,
\end{equation}
where $s_1$ and $s_2$ are the sums of activations before and after pruning. The
premise is that the representation is heavily redundant, so most of it can be
discarded and the sparse surviving support---rescaled to compensate---still
classifies ID inputs while sharpening OOD separation. Because the threshold $t$ is
computed per sample from the whole vector, ASH is whole-vector and not reducible to
an element-wise map, and it uses no training statistics, making it data-free and
therefore white-box (cf.\ Section~\ref{subsec:access}). An extreme ablation, ASH-Rand, replaces the
surviving values with random ones; we return to it in the discussion, as it bears on
what ASH actually exploits.

\subsubsection{Optimal-FS}
\label{subsec:optfs}

Optimal-FS~\citep{optimalfs} recasts element-wise shaping as an optimization
problem. Writing the per-coordinate reshaping as a multiplicative factor
$\theta(z_i) = f(z_i)/z_i$, the shaped maximum logit is
\begin{equation}
\label{eq:optfs}
\bar{\ell}_{\max}(\bar{\mathbf{z}}) = \sum_{i=1}^{M} \theta(z_i)\, w^{\max}_i\, z_i ,
\end{equation}
so the question becomes which feature \emph{value ranges} should be up- or
down-weighted. Optimal-FS partitions the feature range into $K$ intervals, treats
$\theta$ as piecewise constant over them, and collects the per-interval contributions
into a $K$-dimensional interval-specific feature impact (ISFI) vector
$\mathbf{I}(\mathbf{z})$, giving $\bar{\ell}_{\max} = \boldsymbol{\theta}^{\top}
\mathbf{I}(\mathbf{z})$. The shaping vector $\boldsymbol{\theta}$ is then optimized to
separate the ID and OOD maximum logits under a norm constraint; solved with OOD data it
closely matches the element-wise rules above, recovering them as approximations to a
common optimum. Dropping the OOD term gives a closed-form solution from ID data alone (the optimization problem
and its reduction are in Appendix~\ref{app:derivations}),
\begin{equation}
\boldsymbol{\theta}^{\star} = \frac{\sqrt{K}}
{\big\| \mathbb{E}_{\mathrm{ID}}[\mathbf{I}(\mathbf{z})] \big\|}\,
\mathbb{E}_{\mathrm{ID}}\big[\mathbf{I}(\mathbf{z})\big] .
\end{equation}
Nearly all of the benefit comes from $K \leq 10$, indicating the optimal shaping
function has a simple form (the full ISFI derivation is in the appendix). Two
caveats matter for this report. First, the optimization is specific to element-wise
shaping, so it unifies ReAct and VRA-P but explicitly \emph{does not} capture
whole-vector methods such as ASH, which the authors name as an open problem.
Second, like ReAct and VRA, it depends on ID data, making it open-box
(cf.\ Section~\ref{subsec:access}).

\subsection{Between-Layer Smoothness}
\label{subsec:blood}

BLOOD~\citep{blood} reads a different signal from the same region of the network.
Rather than reshaping $\mathbf{z}$, it measures how smoothly the network transforms
its representations, on the premise that a trained model maps ID inputs through
smoother transformations than OOD ones. Smoothness at layer $l$ is the squared
Frobenius norm of that layer's Jacobian $\mathbf{J}_l$,
\begin{equation}
\label{eq:blood}
\phi_l(\mathbf{x}) = \big\| \mathbf{J}_l(\mathbf{h}_l) \big\|_F^{2},
\end{equation}
estimated without forming the Jacobian through an unbiased random-probe estimator
$\hat{\phi}_l$ (Appendix~\ref{app:derivations}). BLOOD uses this
quantity directly as its score, $\hat{\phi}_{L-1}(\mathbf{x})$, with
the \emph{inverted} orientation relative to the shaping cluster: a higher value
indicates OOD. The detector $G_\lambda$ still applies, only with the sides of
$\lambda$ reversed. Note that this score is not expressible in the
shaping composition $s \circ \mathbf{W} \circ f$ of Section~\ref{subsec:shaping}: it
reads the Jacobian and never factors through the head or a shaping map, which is
precisely why BLOOD sits outside feature shaping. The method has two variants,
BLOOD$_M$ averaging $\hat{\phi}_l$ over all layers and BLOOD$_L$ using only the final
transformation, the Jacobian of the head with respect to the penultimate
representation $\mathbf{z}$. The authors focus on BLOOD$_L$, which outperforms
BLOOD$_M$ more often, attributing this to the final layer undergoing the most
training-driven change; this coincides with our penultimate-layer scope, so we take
BLOOD$_L$ as canonical. BLOOD requires only the model's weights and no training data,
making it white-box (cf.\ Section~\ref{subsec:access}).



\section{Experiments}
\label{sec:experiments}

The five methods were never evaluated under a shared protocol, therefore no unified
leaderboard exists: they differ in backbone, dataset suite, and scale. The organizing
axis the evidence does support is \emph{architecture}, not modality: the shaping
cluster is at home on convolutional networks (CNNs), while the last-layer smoothness of
BLOOD$_L$ is at home on transformers, each degrading sharply on the other's
architecture. We therefore read the evidence as two architecture-defined
\emph{islands}---a CNN-optimized shaping island and a transformer-optimized smoothness
island---and treat that split as itself a finding. All figures below are reproduced from
the cited papers, not independently re-run, and a direct cross-island ranking is
precluded by their differing scales and suites.

\subsection{Feature Shaping: The CNN Island}
\label{subsec:exp-vision}

Optimal-FS~\citep{optimalfs} provides the only head-to-head evaluation of the shaping
cluster, re-running ReAct, ASH (P/B/S), VRA-P, and its own method across eight
ImageNet-1k backbones under one protocol. On ResNet-50---the architecture all of these
methods were designed on---the shaping methods are strong and close together (AUROC
$\approx 90.3$ for ASH-S, $88.6$ for VRA-P, $86.5$ for ReAct, $88.6$ for Optimal-FS).
The picture inverts off ResNet-50. On transformer and multi-layer-perceptron (MLP)
backbones the most aggressive shaping rules collapse toward chance: ASH-S falls to
$18.5$ AUROC on ViT-B-16 (a Vision Transformer) and
$33.9$ on MLP-B, and VRA-P behaves similarly, dragging their eight-backbone averages to
$\approx 38$. ReAct degrades far more gracefully (averaging $81.3$, never falling below
the high $70$s), and Optimal-FS holds across the board ($83.0$ average, $82.7$ on
ViT-B-16). So the cluster's apparent unity is conditional on its home architecture; the
clip-based ReAct and the ID-optimized Optimal-FS generalize, while the percentile-prune
methods do not. On the smaller-scale CIFAR benchmarks, feature shaping as a whole
underperforms simple logit baselines, with Optimal-FS competitive but not dominant.

\subsection{Between-Layer Smoothness: The Transformer Island}
\label{subsec:exp-text}

BLOOD~\citep{blood} is evaluated mainly on text classification with RoBERTa and
ELECTRA over eight ID datasets. Among methods requiring no training data
(white/black-box), BLOOD$_L$ is the strongest on most datasets, and more consistently
than the all-layer variant BLOOD$_M$; the authors accordingly focus on BLOOD$_L$,
attributing its edge to the final layer absorbing the most training-driven change. An
image-classification appendix shows this signal is architectural, not textual:
BLOOD$_L$ stays strong on a fine-tuned Vision Transformer but collapses to near-chance
on from-scratch CNNs (ResNet-34/50). BLOOD$_L$ is thus at home on transformers
regardless of modality. Consistent with this, the shaping methods appearing in BLOOD's tables, ASH and
ReAct, are only middling, indicating CNN-tuned shaping does not transfer to
transformers.

\subsection{A Shared Ceiling}
\label{subsec:exp-ceiling}

A distance-based open-box method, Mahalanobis distance~\citep{mahalanobis}, bounds both
islands from above. In BLOOD's comparison it outperforms BLOOD in every setup but one
(ELECTRA on TREC), and it is consistently the top scorer on the transformer
backbones. Its dominance is, however, architecture-specific rather than universal: on
ImageNet-scale ConvNets it is weak (AUROC $\approx 55$ on ResNet-50 in VRA's
table~\citep{vra}). The islands therefore deliver a structural message, not a numerical ranking:
each paradigm wins on its home architecture, and a simple distance detector tops both
wherever its distributional assumptions hold.

\section{Discussion}
\label{sec:discussion}

Placing the five methods together raises three questions: why the same method succeeds
on some architectures and fails on others (Section~\ref{subsec:disc-arch}), how far the
shaping cluster unifies (Section~\ref{subsec:disc-unify}), and whether the OOD signal
lies in the feature values or in the transformation applied to them
(Section~\ref{subsec:disc-carrier}).

\subsection{Impact of Model Architectures}
\label{subsec:disc-arch}

The transfer failures of Section~\ref{subsec:exp-vision} follow from the mechanisms,
not from incidental tuning. Optimal-FS reports the pattern directly: for shaping
methods, convolutional-network performance is negatively correlated with performance on
other architectures~\citep{optimalfs}. ASH's pruning assumes a redundant, sparse
representation: excess features of an over-parameterized network~\citep{ash}, a ReLU
(rectified linear unit) convolutional premise. BLOOD shows the mirror-image dependence:
BLOOD$_L$ reads smoothness in the final transformation, which carries the OOD signal
only where training concentrates there---fine-tuned transformers. Its own
image-classification results bear this out, with BLOOD$_L$ staying strong on a Vision
Transformer but collapsing on from-scratch CNNs, whose learning sits in the lower
layers~\citep{blood}. Feature shaping is thus a CNN technology and BLOOD$_L$ a
transformer one, each encoding the architecture it was built on.

ReAct gives a partial cause. It attributes the abnormal OOD activations to
batch-normalization statistics estimated on ID data and applied to OOD inputs, making
batch normalization the exploited signal rather than an obstacle~\citep{react}. ReAct
also helps under weight and group normalization~\citep{react}, so the effect is not
specific to batch normalization. The architecture dependence is therefore real but its
root cause is unresolved across the cluster.

\subsection{How Far the Shaping Cluster Unifies}
\label{subsec:disc-unify}

Optimal-FS defines shaping as optimization over functions
$f : \mathbb{R}^{M} \to \mathbb{R}^{M}$ and proves that ReAct and VRA-P, rewritten via
$\theta(z) = f(z)/z$ (Eq.~\ref{eq:optfs}), approximate the optimal element-wise
function~\citep{optimalfs}.
This optimum is simple: the large jump above the baseline already appears at a
single interval ($K=1$), which prunes only the $0.1\%$ tails of the feature
distribution, and almost all remaining benefit is captured by $K \leq
10$~\citep{optimalfs}. This explains why ReAct, a single clip, stays competitive
with the optimization-derived method on the architectures both were tuned for:
little structure beyond the tails remains to recover.

The unification is bounded, and Optimal-FS draws the boundary. Its tractable instance
restricts to element-wise $f$ and, in its own statement, fails to capture whole-vector
methods such as ASH, which it still places inside the general
$\mathbb{R}^{M} \to \mathbb{R}^{M}$ class~\citep{optimalfs}. ASH is thus excluded by the
reduction, not the definition, and the whole-vector extension is named as open future
work~\citep{optimalfs}. Extending this boundary to a case Optimal-FS does not treat, we
observe that BLOOD lies outside the class: its score is the Frobenius norm of a
Jacobian (Eq.~\ref{eq:blood})~\citep{blood}, not a map returning a reshaped feature
vector, so no extension of the optimization can absorb it. This gives three tiers: ReAct and VRA-P inside the solved
problem, ASH inside the general class but outside the solved instance, and BLOOD outside
the class. The first gap is one of tractability; the second is categorical.

Whether the whole-vector extension is reachable is open. Two findings bear on it, and
connecting them is our reading rather than a stated result. The element-wise optimum is coarse~\citep{optimalfs}, and ASH's signal is also
coarse: ASH-Rand replaces the surviving activations with random values yet still
beats the logit baselines, locating the signal in the surviving \emph{support}
rather than the value profile, while ASH-B and ASH-S---which discard the profile
but preserve the activation mass (Eq.~\ref{eq:ash})---show the global scale is what that support
must carry~\citep{ash}.

\subsection{Where the Signal Lives}
\label{subsec:disc-carrier}
 
Because BLOOD is not a shaping function, it serves as a counterpoint that differs from
the cluster on one variable: what is measured at the top of the network. The four shaping
methods edit the activation vector; BLOOD leaves it unchanged and measures the smoothness
of the transformation applied to it~\citep{blood}. This poses a question neither paradigm
sets out to answer: is the OOD signal a property of the feature values or of the map
applied to them?

We do not resolve which is fundamental. The juxtaposition does narrow the cluster's
premise. BLOOD shows the activation vector is not the only carrier of OOD signal at the
penultimate layer, and ASH-Rand shows that within the cluster the fine values are not the
operative part. The shared assumption---that the signal can be recovered by reshaping the
penultimate activation vector---is therefore narrower than its framing, though not
incorrect.

\section{Conclusion}
\label{conc}

We read five post-hoc, feature-based OOD detectors (the shaping cluster of ReAct,
ASH, VRA, and Optimal-FS, together with BLOOD) as a single group united by the
penultimate layer they all act on. Three conclusions follow. First, architecture, not
mechanism, decides success: shaping rules tuned on ResNet-50 collapse on transformer
and MLP backbones, and only the conservative clip of ReAct and the ID-optimized shaping
of Optimal-FS transfer, so the cluster's apparent unity is an artifact of its shared
home architecture. Second, the unification is real but bounded: Optimal-FS's
optimization recovers the element-wise rules as approximations to one ID-only optimum,
but whole-vector ASH lies outside the solved instance and BLOOD, scoring a Jacobian
norm, lies outside the shaping class entirely. Third, the OOD signal is narrower than
the shaping premise suggests: BLOOD recovers signal without touching the feature
values, and ASH-Rand recovers it without their fine structure. None of these methods is
best in a setting-free sense; each encodes the architecture it was built
on---feature shaping a CNN technology, BLOOD$_L$ a transformer one---and a simple
distance baseline outperforms both paradigms wherever its assumptions hold. We leave open whether the whole-vector case admits the same optimization account
that unifies the element-wise rules.


\bibliography{ood_base}
\bibliographystyle{abbrvnat}

\newpage
%%%%%%%%%%%%%%%%%%%%%%%%%%%%%%%%%%%%%%%%%%%%%%%%%%%%%%%%%%%%%%%%%%%%%%%%%%%%%%%%%%%%%%%%%%%%%%%%%%%%
%% Use of Large Language Models
%%%%%%%%%%%%%%%%%%%%%%%%%%%%%%%%%%%%%%%%%%%%%%%%%%%%%%%%%%%%%%%%%%%%%%%%%%%%%%%%%%%%%%%%%%%%%%%%%%%%
\section*{Use of Large Language Models}

The LLM Claude Opus 4.8 was used in the initial screening of scientific literature and the explanation of certain concepts during the research phase. During the writing of the report, the model was used to propose changes to certain technical passages, the unification of the mathematical notation, the solution of syntactical errors and sanity checks regarding overall structure of the report. 
%%%%%%%%%%%%%%%%%%%%%%%%%%%%%%%%%%%%%%%%%%%%%%%%%%%%%%%%%%%%%%%%%%%%%%%%%%%%%%%%%%%%%%%%%%%%%%%%%%%%
%% Declaration of Authorship
%%%%%%%%%%%%%%%%%%%%%%%%%%%%%%%%%%%%%%%%%%%%%%%%%%%%%%%%%%%%%%%%%%%%%%%%%%%%%%%%%%%%%%%%%%%%%%%%%%%%




{\parindent 0cm
%%%%%%%%%%%%%%%%%%%%%%%%%%German%%%%%%%%%%%%%%%%%%%%%%%%%%%%%%
\section*{Declaration of Authorship}
Ich erkläre hiermit gemäß § 9 Abs. 12 APO, dass ich die vorstehende Seminararbeit selbstständig verfasst und keine anderen als die angegebenen Quellen und Hilfsmittel benutzt habe. Des Weiteren erkläre ich, dass die digitale Fassung der gedruckten Ausfertigung der Seminararbeit ausnahmslos in Inhalt und Wortlaut entspricht und zur Kenntnis genommen wurde, dass diese digitale Fassung einer durch Software unterstützten, anonymisierten Prüfung auf Plagiate unterzogen werden kann.\\
\vspace{2\baselineskip}
  
Bamberg, \today

\rule[0.5em]{14em}{0.5pt} \hspace{0.25\linewidth}\rule[0.5em]{14em}{0.5pt}
\vspace{1em}
\hspace{4em} (Ort, Datum) \hspace{0.51\linewidth} (Unterschrift)
}



%%%%%%%%%%%%%%%%%%%%%%%%%%%%%%%%%%%%%%%%%%%%%%%%%%%%%%%%%%%%


%%%%%%%%%%%%%%%%%%%%%%%%%%%%%%%%%%%%%%%%%%%%%%%%%%%%%%%%%%%%
\newpage

\appendix


\section{Derivations}
\label{app:derivations}

Writing the per-interval contributions as the ISFI vector $\mathbf{I}(\mathbf{z})$
(Section~\ref{subsec:optfs}), Optimal-FS optimizes the shaping vector
$\boldsymbol{\theta}$ to separate the ID and OOD maximum logits under a norm constraint
that prevents trivial rescaling,
\begin{equation}
\max_{\boldsymbol{\theta}} \;\;
\boldsymbol{\theta}^{\top}
\Big( \mathbb{E}_{\mathrm{ID}}\big[\mathbf{I}(\mathbf{z})\big]
- \mathbb{E}_{\mathrm{OOD}}\big[\mathbf{I}(\mathbf{z})\big] \Big),
\qquad \text{s.t.} \;\; \|\boldsymbol{\theta}\| = \sqrt{K}.
\end{equation}
Dropping the OOD term, on the argument that it contributes little beyond the ID
direction, leaves the closed-form $\boldsymbol{\theta}^{\star}$ of
Section~\ref{subsec:optfs} \cite{optimalfs}.


For BLOOD, the layer-$l$ smoothness $\phi_l = \|\mathbf{J}_l\|_F^{2}$
(Eq.~\ref{eq:blood}) is estimated without forming the Jacobian, using $N$ pairs of
isotropic random probes $\mathbf{v}_i, \mathbf{w}_i$ (zero mean, unit variance),
\begin{equation}
\hat{\phi}_l(\mathbf{x}) = \frac{1}{N} \sum_{i=1}^{N}
\big( \mathbf{w}_i^{\top} \mathbf{J}_l(\mathbf{h}_l)\, \mathbf{v}_i \big)^{2},
\qquad
\mathbb{E}\big[\hat{\phi}_l(\mathbf{x})\big] = \big\| \mathbf{J}_l \big\|_F^{2},
\end{equation}
the isotropy of the probes being what makes the estimator unbiased \cite{blood}.

\newpage
\section{Notation}


% Requires: \usepackage{booktabs} and \usepackage{longtable}

\begin{longtable}{@{}ll@{}}
\caption{Unified notation used throughout this report. \\ 
Symbols are grouped by the work that introduces the corresponding quantity in the OOD-detection
context. Typographic conventions---bold vectors and matrices,
\(\mathds{1}\{\cdot\}\), and \(\top\) for transpose---follow
\citet{goodfellow2016}. Where a symbol is renamed for cross-paper
consistency, the original symbol is given in parentheses.}
\label{tab:notation}\\
\toprule
\textbf{Symbol} & \textbf{Meaning}\\
\midrule
\endfirsthead
\toprule
\textbf{Symbol} & \textbf{Meaning}\\
\midrule
\endhead
\midrule
\multicolumn{2}{r}{\textit{continued on next page}}\\
\endfoot
\bottomrule
\endlastfoot

\multicolumn{2}{@{}l}{\textit{General notation}}\\
\(x\)                     & Input sample\\
\(z\)            & Penultimate-layer feature (activation) vector\\
\(z_i\)                   & \(i\)-th element of \(z\)\\
\(\bar{z}\)      & Feature vector after shaping / rectification\\
\(M\)                     & Penultimate feature dimension\\
\(C\)                     & Number of classes\\
\(\ell_i,\ \ell^{\max}\)  & \(i\)-th logit; max logit\\
\(\mathbf{J}\)            & Jacobian matrix\\
\(\lVert\cdot\rVert_F\)   & Frobenius norm\\
\(\mathds{1}(\cdot)\)   & Indicator (1 if condition holds, else 0)\\
\(\mathbb{E}[\cdot]\)     & Expectation\\
\midrule

\multicolumn{2}{@{}l}{\citet{react}}\\
\(c\)                     & Activation clip threshold, \(\mathrm{ReAct}(z;c)=\min(z,c)\)\\
\(s(\cdot)\)              & OOD scoring function (MSP, MaxLogit, Energy, \dots)\\
\(\lambda\)               & OOD decision threshold\\
\(G_\lambda\)             & OOD detector (in/out rule on \(s\))\\
\midrule

\multicolumn{2}{@{}l}{\citet{vra}}\\
\(p_{\text{in}},\ p_{\text{out}}\) & ID / OOD marginal densities over activations\\
\(g(\cdot),\ g^{*}(\cdot)\)        & Rectification function; variational optimum\\
\(\tau\)                  & Fidelity--gap trade-off coefficient \,(orig.\ \(\lambda\))\\
\(\alpha,\ \beta\)        & Low / high activation thresholds\\
\(\gamma\)                & VRA+ amplification offset on the intermediate band\\
\(\eta_\alpha,\ \eta_\beta\) & ID quantiles defining \(\alpha,\beta\)\\
\midrule

\multicolumn{2}{@{}l}{\citet{optimalfs}}\\
\(\mathbf{w}^{\max}\)     & Weight vector of the predicted (max-logit) class\\
\(\theta(\cdot)\)         & Element-wise shaping function, \(\theta(z):=f(z)/z\)\\
\(\boldsymbol{\theta}\)   & Per-bucket shaping vector, \(\boldsymbol{\theta}\in\mathbb{R}^{K}\)\\
\(I_{[a,b)}(\mathbf{z})\) & Interval-specific feature impact (ISFI) of bucket \([a,b)\)\\
\(a_k,\ b_k\)             & Edges of the \(k\)-th bucket\\
\(K\)                     & Number of buckets\\
\midrule

\multicolumn{2}{@{}l}{\citet{blood}}\\
\(\phi_l(x)\)             & Score: squared Frobenius norm of the layer-\(l\) Jacobian\\
\(h_l\)                   & Layer-\(l\) representation \,(\(h_{L-1}=\mathbf{z}\))\\
\(N\)                     & Number of random probes \,(orig.\ \(M\))\\
\(\mathbf{v}_i,\ \mathbf{w}_i\) & Gaussian probe / projection vectors\\
\(\mathbf{u}_i\)          & Jacobian--vector product, \(\mathbf{u}_i=\mathbf{J}_{L-1}\mathbf{v}_i\)\\
\midrule

\multicolumn{2}{@{}l}{\citet{ash}}\\
\(p\)                     & Pruning percentile\\
\(t\)                     & Per-sample threshold, \(t=\mathrm{Perc}_p(\mathbf{z})\)\\
\(m_i\)                   & Binary keep-mask, \(m_i=\mathds{1}\{z_i\ge t\}\)\\
\(n\)                     & Number of surviving activations\\
\(s_1,\ s_2\)             & Pre- / post-pruning activation sums\\

\end{longtable}

\end{document}
